\documentclass{training-exam}
\schoollogo{../../../resources/latex/crack.pdf}
\usepackage[utf8]{inputenc}
\usepackage[T1]{fontenc}
\usepackage{lmodern}
\setlang{french}
\usepackage{hyperref}
\usepackage{graphicx}
\usepackage{tabularx}
\usepackage{fvextra}
\fvset{breaklines=true,breakanywhere=true,breaksymbol={},breaksymbolleft={},breaksymbolright={},breakanywheresymbolpre={},breakanywheresymbolpost={}}
\makeatletter
\g@addto@macro\UrlBreaks{%
  \do\a\do\b\do\c\do\d\do\e\do\f\do\g\do\h\do\i\do\j\do\k\do\l\do\m%
  \do\n\do\o\do\p\do\q\do\r\do\s\do\t\do\u\do\v\do\w\do\x\do\y\do\z%
  \do\A\do\B\do\C\do\D\do\E\do\F\do\G\do\H\do\I\do\J\do\K\do\L\do\M%
  \do\N\do\O\do\P\do\Q\do\R\do\S\do\T\do\U\do\V\do\W\do\X\do\Y\do\Z%
  \do\0\do\1\do\2\do\3\do\4\do\5\do\6\do\7\do\8\do\9%
  \do\<\do\>\do\(\do\)\do\'\do\/}
\schoollogo{../../../resources/latex/crack.pdf}
\makeatother


\begin{document}

\title{Unseen Patch}
\subtitle{Analyse d'un correctif inconnu}
\member{Ted ``Tedsig42'' Kouhouenou}{ted.sig42@tutamail.com}
\summary{Examen de fin de piscine. Vous recevez un correctif que vous n'avez jamais vu. Vous devez identifier sa famille, formuler son hypothèse, trouver un twin, le vérifier et le trier.}
\examduration{3 hours}
\examscore{20}
\maketitle

\startexam

\chapter{Règles}

\begin{itemize}
  \item Durée : 3 heures. Aucun accès Internet, aucune note, aucun outil
        d'IA.
  \item Le correctif et les sources nécessaires sont fournis localement.
  \item Vous pouvez utiliser \texttt{git}, \texttt{grep}, \texttt{find},
        \texttt{diff} et un éditeur de texte. \texttt{ctags} est optionnel.
  \item Aucun exploit, aucun test sur une cible réelle.
  \item Les réponses non justifiées ne rapportent pas de points.
\end{itemize}

\examlevel{Cadrage}

On vous fournit un commit de correction, son message, et le fichier
concerné. Sans lire le code plus de dix minutes, répondez :

\begin{enumerate}
  \item Quel est le composant ? Quelle famille principale ?
  \item Quelle ligne du diff porte l'hypothèse ?
  \item Quel chemin le correctif semble-t-il protéger ?
  \item Quel chemin n'est pas mentionné dans le message de commit ?
\end{enumerate}

\exampoints{4}

\examlevel{Extraction}

Extrayez les lignes ajoutées et supprimées. Classez-les en trois
catégories : test, action, commentaire.

Puis répondez :

\begin{enumerate}
  \item La correction ajoute-t-elle une vérification, une initialisation,
        ou un changement de type ?
  \item La vérification porte-t-elle sur la valeur, l'index, la taille ou
        l'état ?
  \item Le correctif déplace-t-il une responsabilité d'une fonction à une
        autre ?
\end{enumerate}

\exampoints{4}

\examlevel{Twins}

Trouvez au moins trois fonctions voisines qui font un travail analogue.
Pour chacune, indiquez :

\begin{enumerate}
  \item le nom et la ligne ;
  \item l'invariant partagé ;
  \item le test présent ou absent ;
  \item le statut : candidat sérieux, faux positif, ou à vérifier.
\end{enumerate}

Vous devez fournir les commandes \texttt{grep} (et \texttt{find} si besoin)
utilisées. \texttt{ctags} n'est pas exigé.

\exampoints{5}

\examlevel{Vérification}

Choisissez le candidat le plus sérieux. Expliquez :

\begin{enumerate}
  \item pourquoi il partage réellement l'hypothèse du correctif ;
  \item pourquoi il n'est pas protégé par un autre invariant ;
  \item quel test permettrait de confirmer ou d'infirmer le candidat sans
        exploit ;
  \item quel serait l'impact si l'hypothèse était fausse.
\end{enumerate}

\exampoints{4}

\examlevel{Triage et conclusion}

Donnez un score de triage au candidat selon la grille :

\begin{center}
\begin{tabularx}{\linewidth}{|l|X|c|}
\hline
\textbf{Critère} & \textbf{Question} & \textbf{Points} \\ \hline
Confiance & Le correctif est-il réellement lié au candidat ? & 0--3 \\ \hline
Proximité & Même chemin ou twin direct ? & 0--3 \\ \hline
Impact & Que se passe-t-il si l'hypothèse est fausse ? & 0--2 \\ \hline
Testabilit\'e & Vérifiable sans exploit ? & 0--2 \\ \hline
\end{tabularx}
\end{center}

Concluez en dix lignes : \emph{ignorer}, \emph{investiguer}, ou
\emph{prioriser}. Justifiez.

\exampoints{3}

\chapter{Rendu}

Un seul fichier \texttt{exam.md} par étudiant, avec les sections suivantes :

\begin{console}
# Cadrage
# Extraction
# Twins
# Verification
# Triage
# Sources
\end{console}

\chapter{Ressources de révision}

\begin{itemize}
  \item Papier, sections 3, 4 et 18.3 : famille, hypothèse, twins.
  \item Dépôt compagnon : \texttt{rules/TRIAGE.md},
        \texttt{validation/methodology.md}.
  \item Optionnel : Universal Ctags : \url{https://docs.ctags.io/}
  \item NVD : \url{https://nvd.nist.gov/}
  \item À réviser : \texttt{git show}, \texttt{grep '\string^+'},
        \texttt{grep -RIn}, \texttt{nft\_is\_active\_next}, \texttt{regsafe},
        \texttt{fdget}.
\end{itemize}

\chapter{Barème}

\begin{center}
\begin{tabularx}{\linewidth}{|l|c|}
\hline
\textbf{Partie} & \textbf{Points} \\ \hline
Cadrage & 4 \\ \hline
Extraction & 4 \\ \hline
Twins & 5 \\ \hline
Vérification & 4 \\ \hline
Triage et conclusion & 3 \\ \hline
& \textbf{20} \\ \hline
\end{tabularx}
\end{center}

\end{document}
